\section{Banco de questões — História e Geografia do Maranhão}

\subsection*{N1 — Fundamentos em nível de concurso}
\begin{questao}{\nivel{N1} 08.01 — Cebraspe (C/E)}A França Equinocial, instalada em 1612, relaciona-se diretamente à fundação de São Luís e não deve ser confundida com a França Antártica.\end{questao}
\begin{questao}{\nivel{N1} 08.02 — Cebraspe (C/E)}A Batalha de Guaxenduba ocorreu no contexto da reação luso-portuguesa à presença francesa no Maranhão.\end{questao}
\begin{questao}{\nivel{N1} 08.03 — FGV}A ocupação holandesa de São Luís ocorreu:\begin{enumerate}[label=\Alph*)]\item antes da França Equinocial.\item no século XVII, posteriormente ao conflito luso-francês.\item durante a Balaiada.\item após a Revolução de 1930.\item no Segundo Reinado.\end{enumerate}\end{questao}
\begin{questao}{\nivel{N1} 08.04 — Cebraspe (C/E)}A Revolta de Beckman esteve associada à insatisfação com mecanismos econômicos coloniais, incluindo a atuação da Companhia de Comércio do Maranhão.\end{questao}
\begin{questao}{\nivel{N1} 08.05 — Cebraspe (C/E)}A adesão do Maranhão à Independência ocorreu de modo imediato em 7 de setembro de 1822, sem resistência ou conflito regional posterior.\end{questao}
\begin{questao}{\nivel{N1} 08.06 — FCC}A Balaiada caracterizou-se por:\begin{enumerate}[label=\Alph*)]\item movimento exclusivamente militar sem participação popular.\item revolta social do Período Regencial com forte participação popular.\item invasão estrangeira.\item rebelião republicana de 1889.\item movimento exclusivamente urbano de São Luís em 1951.\end{enumerate}\end{questao}
\begin{questao}{\nivel{N1} 08.07 — Cebraspe (C/E)}O Maranhão limita-se com Piauí, Tocantins e Pará, além de possuir litoral voltado para o oceano Atlântico.\end{questao}
\begin{questao}{\nivel{N1} 08.08 — Cebraspe (C/E)}Parnaíba, Gurupi e Tocantins-Araguaia aparecem no edital como bacias limítrofes, enquanto Mearim e Itapecuru são exemplos relevantes de bacias maranhenses.\end{questao}
\begin{questao}{\nivel{N1} 08.09 — Vunesp}No sistema de unidades de conservação, uma APA é, em regra:\begin{enumerate}[label=\Alph*)]\item unidade de proteção integral necessariamente sem ocupação humana.\item unidade de uso sustentável.\item área sem qualquer regime de proteção.\item sinônimo de parque nacional.\item área destinada exclusivamente à mineração.\end{enumerate}\end{questao}
\begin{questao}{\nivel{N1} 08.10 — Cebraspe (C/E)}A Mata dos Cocais é formação de transição na qual se destacam palmeiras como o babaçu.\end{questao}

\subsection*{N2 — Aplicação}
\begin{questao}{\nivel{N2} 08.11 — FGV}Uma questão associa Daniel de La Touche, fundação de São Luís e França Equinocial. A relação é:\begin{enumerate}[label=\Alph*)]\item historicamente coerente.\item incorreta, pois La Touche liderou a Balaiada.\item incorreta, pois São Luís foi fundada por holandeses.\item incorreta, pois França Equinocial ocorreu no século XIX.\item incorreta, pois o episódio se deu no Piauí.\end{enumerate}\end{questao}
\begin{questao}{\nivel{N2} 08.12 — Cebraspe (C/E)}A Batalha do Jenipapo, embora travada no Piauí, integra o contexto regional de consolidação da Independência no norte do Brasil e se relaciona à adesão do Maranhão ao Império.\end{questao}
\begin{questao}{\nivel{N2} 08.13 — Cebraspe (C/E)}Reduzir a Revolta de Beckman a movimento de independência nacional seria historicamente inadequado.\end{questao}
\begin{questao}{\nivel{N2} 08.14 — FCC}O Vitorinismo está associado principalmente:\begin{enumerate}[label=\Alph*)]\item à presença francesa no século XVII.\item à influência política de Vitorino Freire no Maranhão do século XX.\item à Guerra do Paraguai.\item à expansão holandesa.\item à criação do Estado do Grão-Pará no século XVII.\end{enumerate}\end{questao}
\begin{questao}{\nivel{N2} 08.15 — Cebraspe (C/E)}A Greve de 1951 deve ser compreendida no contexto de tensões políticas e mobilização popular do Maranhão do pós-guerra.\end{questao}
\begin{questao}{\nivel{N2} 08.16 — FGV}A diversidade ambiental do Maranhão é melhor explicada por sua posição:\begin{enumerate}[label=\Alph*)]\item exclusivamente amazônica.\item exclusivamente sertaneja.\item de transição entre domínios amazônicos, cerrados, cocais e ambientes costeiros.\item inteiramente inserida na Caatinga.\item sem contato com o oceano.\end{enumerate}\end{questao}
\begin{questao}{\nivel{N2} 08.17 — Cebraspe (C/E)}A pluviosidade tende a apresentar diferenças espaciais no Maranhão, influenciando regime fluvial, vegetação e uso agropecuário.\end{questao}
\begin{questao}{\nivel{N2} 08.18 — Cebraspe (C/E)}Mearim e Itapecuru possuem relevância para ocupação, abastecimento e atividades econômicas, razão pela qual sua leitura não deve se limitar à memorização do nome da bacia.\end{questao}
\begin{questao}{\nivel{N2} 08.19 — FGV}A Baixada Maranhense relaciona-se de modo marcante a:\begin{enumerate}[label=\Alph*)]\item dinâmica sazonal de inundação.\item relevo exclusivamente montanhoso.\item ausência de corpos d'água.\item clima desértico.\item mineração de carvão como única atividade.\end{enumerate}\end{questao}
\begin{questao}{\nivel{N2} 08.20 — Cebraspe (C/E)}A expansão da soja no sul e no leste do Maranhão integra a dinâmica recente do agronegócio associada ao MATOPIBA.\end{questao}
\begin{questao}{\nivel{N2} 08.21 — FCC}O Porto do Itaqui e os corredores ferroviários são relevantes porque:\begin{enumerate}[label=\Alph*)]\item reduzem a integração do estado com mercados externos.\item articulam fluxos de exportação, combustíveis, minérios e produtos agrícolas.\item possuem função apenas turística.\item substituem toda a malha rodoviária.\item se localizam fora do Maranhão.\end{enumerate}\end{questao}
\begin{questao}{\nivel{N2} 08.22 — Cebraspe (C/E)}A concentração de serviços e funções administrativas em São Luís ajuda a explicar sua centralidade urbana no estado.\end{questao}
\begin{questao}{\nivel{N2} 08.23 — Cebraspe (C/E)}Imperatriz exerce importante centralidade regional e articula fluxos com estados vizinhos, especialmente na porção sudoeste do Maranhão.\end{questao}
\begin{questao}{\nivel{N2} 08.24 — FGV}Babaçu, bumba-meu-boi e Porto do Itaqui pertencem, respectivamente, aos eixos:\begin{enumerate}[label=\Alph*)]\item extrativismo vegetal, cultura e logística.\item mineração, pecuária e religião.\item agricultura irrigada, indústria automobilística e pesca.\item cultura, hidrografia e climatologia.\item política, relevo e demografia.\end{enumerate}\end{questao}
\begin{questao}{\nivel{N2} 08.25 — Cebraspe (C/E)}Urbanização e crescimento populacional não implicam, por si, distribuição homogênea de infraestrutura e serviços públicos.\end{questao}

\subsection*{N3 — Nível de prova}
\begin{questao}{\nivel{N3} 08.26 — Cebraspe (C/E)}Associar a Batalha de Guaxenduba à expulsão dos holandeses é incorreto, pois o episódio integra o conflito contra os franceses.\end{questao}
\begin{questao}{\nivel{N3} 08.27 — FGV}A melhor sequência cronológica é:\begin{enumerate}[label=\Alph*)]\item Balaiada -- França Equinocial -- Beckman -- Independência.\item França Equinocial -- Guaxenduba -- ocupação holandesa -- Beckman.\item Beckman -- Guaxenduba -- Balaiada -- França Equinocial.\item Independência -- França Equinocial -- holandeses -- Balaiada.\item 1930 -- Balaiada -- Beckman -- Guaxenduba.\end{enumerate}\end{questao}
\begin{questao}{\nivel{N3} 08.28 — Cebraspe (C/E)}A Balaiada combina conflitos entre elites locais e forte participação de grupos populares, não sendo adequadamente explicada por uma única causa.\end{questao}
\begin{questao}{\nivel{N3} 08.29 — Cebraspe (C/E)}A Revolução de 1930 alterou relações políticas entre poder central e elites estaduais, afetando também a organização política do Maranhão.\end{questao}
\begin{questao}{\nivel{N3} 08.30 — FGV}Considere: I. França Equinocial; II. Revolta de Beckman; III. Balaiada. Os três eventos pertencem, respectivamente, aos contextos:\begin{enumerate}[label=\Alph*)]\item disputa colonial; crise do sistema colonial; Período Regencial.\item República; Império; Colônia.\item Império; República; Colônia.\item Estado Novo; Colônia; República.\item todos ao século XX.\end{enumerate}\end{questao}
\begin{questao}{\nivel{N3} 08.31 — Cebraspe (C/E)}A posição do Maranhão entre Amazônia e Nordeste contribui para mosaico de vegetação e para diferenças internas de precipitação e uso da terra.\end{questao}
\begin{questao}{\nivel{N3} 08.32 — FGV}Se uma área apresenta babaçuais extensos em zona de transição, a formação vegetal mais diretamente associada é:\begin{enumerate}[label=\Alph*)]\item Mata dos Cocais.\item tundra.\item pradaria pampeana.\item araucária.\item pantanal.\end{enumerate}\end{questao}
\begin{questao}{\nivel{N3} 08.33 — Cebraspe (C/E)}Os manguezais do litoral maranhense possuem relação com ambientes estuarinos e costeiros e não podem ser classificados simplesmente como cerrado.\end{questao}
\begin{questao}{\nivel{N3} 08.34 — FCC}A diferença fundamental entre APA e parque nacional reside em que:\begin{enumerate}[label=\Alph*)]\item ambos são idênticos.\item APA é categoria de uso sustentável, enquanto parque nacional integra proteção integral.\item parque nacional é sempre privado.\item APA não possui objetivo ambiental.\item parque nacional admite qualquer exploração mineral sem restrição.\end{enumerate}\end{questao}
\begin{questao}{\nivel{N3} 08.35 — Cebraspe (C/E)}Uma política de uso do solo no Maranhão deve considerar que relevo, drenagem e regime de chuvas influenciam erosão, inundação e aptidão agrícola.\end{questao}
\begin{questao}{\nivel{N3} 08.36 — FGV}A Estrada de Ferro Carajás possui importância para o Maranhão porque:\begin{enumerate}[label=\Alph*)]\item articula corredor mineral e logístico até a região portuária de São Luís.\item conecta exclusivamente cidades do litoral piauiense.\item é ferrovia urbana de passageiros da capital.\item substitui o Porto do Itaqui.\item não possui relação com exportações.\end{enumerate}\end{questao}
\begin{questao}{\nivel{N3} 08.37 — Cebraspe (C/E)}A economia maranhense combina produção familiar, agronegócio de grãos, pecuária, extrativismo, indústria e serviços, de modo que classificá-la por um único setor é reducionista.\end{questao}
\begin{questao}{\nivel{N3} 08.38 — Cebraspe (C/E)}A expansão de grandes corredores logísticos pode gerar crescimento econômico sem eliminar automaticamente desigualdades territoriais e sociais.\end{questao}
\begin{questao}{\nivel{N3} 08.39 — FGV}Em análise regional, São Luís e Imperatriz devem ser vistas como:\begin{enumerate}[label=\Alph*)]\item centros urbanos com funções e áreas de influência distintas.\item municípios sem função regional.\item áreas rurais exclusivamente agrícolas.\item portos marítimos equivalentes.\item cidades sem conexão rodoviária.\end{enumerate}\end{questao}
\begin{questao}{\nivel{N3} 08.40 — Cebraspe (C/E)}Bumba-meu-boi e tambor de crioula expressam a formação cultural plural do Maranhão e não devem ser reduzidos a manifestações de uma única matriz étnica.\end{questao}

\subsection*{N4 — Avançado e integração}
\begin{questao}{\nivel{N4} 08.41 — FGV}Uma questão relaciona clima mais sazonal no sul do Maranhão, expansão de grãos e integração por corredores logísticos. A relação é plausível porque:\begin{enumerate}[label=\Alph*)]\item condições naturais e infraestrutura influenciam localização das atividades econômicas.\item agricultura independe de clima e logística.\item o sul do estado é integralmente manguezal.\item grãos são produzidos apenas na ilha de São Luís.\item logística não interfere na competitividade agrícola.\end{enumerate}\end{questao}
\begin{questao}{\nivel{N4} 08.42 — Cebraspe (C/E)}A leitura integrada do território permite relacionar desmatamento, expansão agropecuária, disponibilidade hídrica e mudanças no uso da terra, temas que não se esgotam na descrição física do relevo.\end{questao}
\begin{questao}{\nivel{N4} 08.43 — FGV}Um município da Baixada apresenta cheias sazonais e ocupação urbana de áreas baixas. O planejamento deve priorizar:\begin{enumerate}[label=\Alph*)]\item ignorar hidrologia local.\item mapear risco, drenagem, ocupação e dinâmica sazonal das águas.\item tratar toda inundação como evento imprevisível.\item eliminar qualquer área úmida.\item usar apenas dados eleitorais.\end{enumerate}\end{questao}
\begin{questao}{\nivel{N4} 08.44 — Cebraspe (C/E)}O fato de o Porto do Itaqui ampliar capacidade de escoamento não implica que benefícios econômicos e infraestrutura social se distribuam automaticamente de forma uniforme no território.\end{questao}
\begin{questao}{\nivel{N4} 08.45 — FGV}A sequência “França Equinocial -- Beckman -- adesão à Independência -- Balaiada -- Vitorinismo” demonstra:\begin{enumerate}[label=\Alph*)]\item transição entre Colônia, Império e República.\item eventos todos do mesmo século.\item episódios exclusivamente militares.\item fatos sem relação com o Maranhão.\item apenas movimentos de independência.\end{enumerate}\end{questao}
\begin{questao}{\nivel{N4} 08.46 — Cebraspe (C/E)}A Batalha do Jenipapo pode ser cobrada em História do Maranhão porque o processo de Independência no norte do Brasil ultrapassou fronteiras provinciais atuais e afetou a adesão maranhense.\end{questao}
\begin{questao}{\nivel{N4} 08.47 — FGV}Ao estudar cultura maranhense em perspectiva geográfica, é correto relacioná-la:\begin{enumerate}[label=\Alph*)]\item à formação histórica, às matrizes populacionais e aos territórios de sociabilidade.\item apenas a entretenimento sem dimensão histórica.\item exclusivamente ao clima.\item apenas a indicadores econômicos.\item somente à legislação federal.\end{enumerate}\end{questao}
\begin{questao}{\nivel{N4} 08.48 — Cebraspe (C/E)}A distinção entre bacias limítrofes e bacias maranhenses ajuda a compreender fronteiras, drenagem interna e integração interestadual.\end{questao}
\begin{questao}{\nivel{N4} 08.49 — FGV}Um estudo sobre desenvolvimento sustentável no Maranhão deveria integrar, no mínimo:\begin{enumerate}[label=\Alph*)]\item ambiente, população, logística e estrutura produtiva.\item apenas o número de municípios.\item exclusivamente relevo.\item somente produção agrícola.\item apenas dados históricos do século XVII.\end{enumerate}\end{questao}
\begin{questao}{\nivel{N4} 08.50 — Cebraspe (C/E)}Compreender o Maranhão exige articular permanências históricas, desigualdades regionais, diversidade ambiental e inserção em cadeias logísticas nacionais e internacionais.\end{questao}
